\subsection{Principal fiber bundles with base B × I}

Let G be a topological group. We shall see that Proposition 2 and its corollaries remain valid when, in each statement, “locally trivial fiber space” is replaced by “principal fiber space with group G”.

Lemma 4. --- Let B be a topological space, let G be a topological group, and let E, E′ be principal fiber bundles with group G and base B. Denote by F the topological space G equipped with the left action of G×G given by \((g,g')\cdot f=g'fg^{-1}\), and denote by \(M = (\mathrm{E}\times_{\mathrm{B}}\mathrm{E}')\times^{\mathrm{G}\times\mathrm{G}}\mathrm{F}\) the locally trivial fiber bundle with fiber type F and base B associated.

The sheaf \(\mathcal{S}\) on B of sections of M is isomorphic to the sheaf on B of isomorphisms of principal fiber bundles from E to E'.

Denote by \(p\), \(p'\) and \(q\) the projections of the B-spaces E, E' and M. For \((x, x') \in \mathrm{E} \times_{\mathrm{B}} \mathrm{E}'\) and \(f \in \mathrm{F}\), denote by \([x, x', f]\) the class in M of the element \(((x, x'), f) \in (\mathrm{E} \times_{\mathrm{B}} \mathrm{E}') \times \mathrm{F}\). Denote by \(\mathcal{M}\) the sheaf on B of isomorphisms of principal fiber bundles from E to E'; its sections over an open subset U of B are the isomorphisms of principal fiber bundles from \(E_{U}\) to \(E_{U}'\).

Denote by \(e\) the identity element of G. Let U be an open subset of B and let \(\varphi\colon\mathrm{E}_{\mathrm{U}}\to\mathrm{E}_{\mathrm{U}}^{\prime}\) be an isomorphism of principal fiber bundles with group G and base U. The map from \(\mathrm{E}_{\mathrm{U}}\) to \((\mathrm{E}\times_{\mathrm{B}}\mathrm{E}^{\prime})\times^{\mathrm{G}\times\mathrm{G}}\mathrm{F}\) defined by \(x\mapsto[x,\varphi(x),e]\) is continuous. For \(g\in\mathrm{G}\) and \(x\in\mathrm{E}_{\mathrm{U}}\), it sends \(x\cdot g\) to \([x\cdot g,\varphi(x\cdot g),e]=[x,\varphi(x),geg^{-1}]=[x,\varphi(x),e]\). There therefore exists a unique continuous map \(\alpha_{\mathrm{U}}(\varphi)\colon\mathrm{U}\to\mathrm{M}\) such that \(\alpha_{\mathrm{U}}(\varphi)(p(x))=[x,\varphi(x),e]\) for all \(x\in\mathrm{E}_{\mathrm{U}}\); we have \(\alpha_{\mathrm{U}}(\varphi)\in\mathcal{S}(\mathrm{U})\).

It is immediate that the maps \(\alpha_{U}\) define a morphism of sheaves \(\alpha\) from M into S.

Lemma 5. --- The morphism of sheaves α is an isomorphism.

Suppose first that the principal fiber bundles E and E' are both trivializable; let us choose sections \(i: B \to E\) and \(i': B \to E'\). There then exists a unique continuous map \(\theta\) from M into \(B \times G\) such that

\[
\theta ([ i (b) \cdot g, i ^ {\prime} (b) \cdot g ^ {\prime}, f ]) = (b, g ^ {\prime} f g ^ {- 1})
\]

continuous for \(b \in \mathrm{B}\), \(g \in \mathrm{G}\), \(g' \in \mathrm{G}\) and \(f \in \mathrm{F}\); it is an isomorphism of B-spaces. Let \(\varphi\) be an isomorphism of principal fiber bundles from E to \(\mathrm{E}'\); there exists a unique continuous map \(\gamma: \mathrm{B} \to \mathrm{G}\) such that \(\varphi(i(b)) = i'(b) \cdot \gamma(b)\) for \(b \in \mathrm{B}\). The image of \(\varphi\) under the map \(\alpha_{\mathrm{B}}\) is the map \(b \mapsto \theta^{-1}(b, \gamma(b))\) from B to M. This implies that \(\alpha_{\mathrm{B}}\) is a bijection.

Consequently, \(\alpha_{\mathrm{U}}\) is a bijection for every open set U of B such that the principal fiber bundles \(\mathrm{E_U}\) and \(\mathrm{E_U}'\) are trivializable. By corollary 2 of I, p. 55, this implies that \(\alpha\) is an isomorphism of sheaves.

PROPOSITION 3. --- Let G be a topological group, let B be a paracompact topological space, and let (E, p) be a principal fiber bundle with group G and base B × I. Set \(E_{1} = \overline{p}^{-1}(B \times \{1\})\) and let \(p_{1}: E_{1} \to B\) be the map \(pr_{1} \circ p \mid E_{1}\). Then (E, p) and \((E_{1} \times I, p_{1} \times Id_{I})\) are isomorphic principal fiber bundles with group G and base B × I.

Let F be the topological space G equipped with the left action of the group \(G \times G\) given by \((g, g') \cdot f = g' fg^{-1}\). Let \((M, q)\) be the locally trivial fiber bundle with base \(B \times I\) and fiber type F associated with the principal fiber bundle \(E \times_{B \times I}(E_1 \times I)\) of the group \(G \times G\). Set \(M_1 = \overline{q}^{-1}(B \times \{1\})\)

and \(q_1 = \mathrm{pr}_1 \circ q \mid M_1\); the B-space \((M_1, q_1)\) is identified with the locally trivial fiber bundle with fiber type F associated with \(E_1 \times_\mathrm{B} E_1\). By Lemma 4, where one takes the principal fiber bundles E and \(E'\) equal to \(E_1\), the B-space \(M_1\) has a section. Since the \(B \times I\)-spaces \((M, q)\) and \((M_1 \times I, q_1 \times \mathrm{Id}_{I})\) are isomorphic (IV, p. 440, Prop. 2), the \(B \times I\)-space \((M, q)\) has a section, which implies that the principal G-bundles E and \(E_1 \times I\) are isomorphic.

COROLLARY 1. --- Let B be a paracompact topological space, let G be a topological group, and let (E,p) be a principal G-bundle over B × I. For \(t \in I\), denote by \((E_{t}, p_{t})\) the principal B-fiber bundle \(i_{t}^{*}E\), where \(i_{t}: B \to B \times \{t\}\) is the map \(b \mapsto (b, t)\). The principal B-fiber bundles \(E_{0}\) and \(E_{t}\) are isomorphic for every \(t \in I\).

COROLLARY 2. --- Let B and B' be topological spaces, let G be a topological group, let E be a principal B-fiber bundle with group G. Let \(f_0\) and \(f_1\) be continuous maps from B' into B. Suppose that the space B' is paracompact. If the maps \(f_0\) and \(f_1\) are homotopic, the principal B'-fiber bundles \(f_0^*E\) and \(f_1^*E\) are isomorphic.

COROLLARY 3. --- Let B be a paracompact topological space and let G be a topological group. If B is homotopy equivalent to a point, then every principal G-bundle over B is trivializable.

Remark. --- An alternative proof of these results would consist in verifying that the isomorphisms of fiber spaces constructed in Prop. 2 and its corollaries are isomorphisms of principal fiber spaces.

